% Independently audited by root, 2026-11-19 15:50 UTC.
% The checked copy is included in part_A5.tex.
\begin{saresult}{A5.013}{Proved under stated hypotheses}{The maximal triviality locus for targets of dimension at most two}
Let $M$ be an o-minimal expansion of a real closed field and let
$f:X\to Y$ be a continuous definably proper definable map.  Suppose
$\dim Y\leq2$.  Then the largest open subset of $Y$ which is a
definable manifold without boundary and over which $f$ is definably
locally trivial is definable.  Thus the remaining maximal-locus
question requires a target of dimension at least three and
positive-dimensional fibres.
\end{saresult}
\begin{proof}
All neighbourhoods and maps used below are definable; parameters
are allowed.  ``Compact'' in this proof means closed and bounded,
or equivalently definably compact, and never asserts ordinary
topological compactness over a non-Archimedean field.

\medskip\noindent\textbf{1. A finite-format consequence of logical compactness.}
We first give the precise uniformity principle used twice below.
Fix the structure $M$, parameters defining a map $f$, and a
definable set $A$ of target points.  A \emph{certificate format}
is a finite list of formulas for sets and graphs, together with
fixed finite combinatorial data.  Its free tuple $b$ supplies all
geometric parameters.  A predicate $C(a,b)$ is first-order if it
requires the indicated graphs to be continuous maps or
homeomorphisms on the indicated domains, requires specified
diagram identities, and requires the domains to be open, closed,
bounded, or to contain specified points or intervals.  For
example, continuity of a graph $\Gamma_h$ at $x$ is the formula
\[
 \forall\epsilon>0\ \exists\delta>0\ \forall x',z,z'
 \bigl((x,z),(x',z')\in\Gamma_h\ \wedge\ |x-x'|<\delta
       \Longrightarrow |z-z'|<\epsilon\bigr).
\]
Single-valuedness, surjectivity, injectivity and inverse continuity
are likewise first-order.  A fixed finite polyhedron can be
represented by its vertex coordinates and a fixed list of faces;
the required incidence and affine-independence conditions are
finite first-order conditions.  A finite piecewise affine map,
possibly discontinuous on the boundaries of its pieces, has the
same kind of description.

Suppose that in every elementary extension $M'\succeq M$, every
$a\in A(M')$ has a certificate of one of a set of formats
$\{C_i(a,b_i):i\in I\}$.  Then finitely many formats suffice.
Indeed, if no finite list sufficed, the elementary diagram of $M$
together with a new constant $a\in A$ and the formulas
$\neg\exists b_i\,C_i(a,b_i)$ for all $i\in I$ would be finitely
satisfiable in $M$.  The compactness theorem would give an
elementary extension containing an $a$ with no certificate,
contrary to the hypothesis.  The compactness theorem used here is
Theorem~2.4, p.17 of Marker's original Orsay notes
\cite{A5-Marker2010}; this paragraph supplies the entire
application, including the elementary-diagram argument.

Every elementary extension here is still o-minimal.  To see why
this use is legitimate, apply cell decomposition to each fixed
one-variable formula with parameters in $M$.  It gives a uniform
finite bound on the number of interval and point components.
The assertion that every parameter fibre is a union of at most
that many intervals and points is a first-order sentence, using
their endpoints.  It holds in the extension, for each formula.
All closed-boundedness and continuity hypotheses used below also
transfer by their first-order formulations.

Here is the particular consequence concerning local PL models.
Assume that in every elementary extension and at every $a\in A$
there is a local isomorphism of the full map over a target
neighbourhood to a finite simplicial map
$p:|K|\to|L|$.  More explicitly, the certificate consists of
open sets $V\subset Y$, $W\subset|L|$, homeomorphisms
\[
 \beta:W\xrightarrow{\sim}V,
 \qquad
 \alpha:p^{-1}(W)\xrightarrow{\sim}f^{-1}(V),
 \qquad f\alpha=\beta p,
 \qquad a\in V.
 \tag{A5.13.1}
\]
The finite complexes and their vertex map are part of the format;
use their canonical rational realizations, so there are no
varying combinatorial types within one format.  All requirements
in (A5.13.1) are first-order.  The preceding argument therefore
gives finitely many such chart formats covering $A$.

For each of their fixed simplicial maps $p_i$, let $G_i\subset|L_i|$
be its manifold-and-local-triviality locus in the original structure
$M$.  No invariance of that unknown goodness property under
elementary extensions is asserted or needed.  By A5.008,
$G_i$ is a union of finitely many open simplices, and hence is
semilinear.  No decision procedure for choosing these finitely
many simplices is required.  In any certificate of the $i$th
format, local triviality and manifoldhood at $a$ are equivalent
to membership of $\beta^{-1}(a)$ in this same $G_i$.  The formula
\[
 a\in A\ \wedge\
 \bigvee_i\exists b_i\,\exists w\,
 \bigl(C_i(a,b_i)\ \wedge\ (w,a)\in\Gamma_{\beta_i,b_i}
                  \ \wedge\ w\in G_i\bigr)
 \tag{A5.13.2}
\]
thus defines the desired locus on $A$.  This argument does not
assume that the unknown good locus has uniform chart complexity:
it only uses universal existence of PL models on the already
definable set $A$.

\medskip\noindent\textbf{2. A semilinear agreement gives an actual local PL model.}
Let $P$ be a compact finite polyhedron, let $F:P\to Z\subset M^m$
be continuous definable, and let $g:P\to M^m$ be a finite
piecewise affine map, not assumed continuous everywhere.  Suppose
$V\subset Z$ is open and
\[
             F(x)=g(x)\qquad(x\in F^{-1}(V)).
 \tag{A5.13.3}
\]
At every $y\in V$ there is a local finite PL model of the full
map $F$, provided the target has a PL chart near $y$.

Write $O=F^{-1}(V)$ and $E=F^{-1}(y)$.  The equality
$E=O\cap g^{-1}(y)$ shows that $E$ is open in the semilinear set
$g^{-1}(y)$; it is also closed there because $E$ is closed in $P$.
It is a definable clopen subset of that semilinear set.  Each
semilinear definably connected component is semilinearly path connected,
so a definable clopen subset is a union of such components.
There are finitely many components.  Consequently $E$ is a
compact semilinear subpolyhedron of $P$.

If $E$ is empty, compactness of $P$ makes $F(P)$ closed and the
empty map over a small neighbourhood is a finite PL model.
Otherwise triangulate $P$ compatibly with $E$, making the latter
a full subcomplex after barycentric subdivision.  The PL function
$\rho:P\to[0,1]$ which is zero on vertices of $E$ and one on the
other vertices has zero set exactly $E$.  On the compact set
$P\setminus O$ its positive values have a positive minimum when
that set is nonempty.  Choose $c>0$ below that minimum (and any
$c>0$ if $P\setminus O$ is empty), and put
\[
                     Q=\{x\in P:\rho(x)\leq c\}.
\]
Then $Q$ is a compact polyhedral neighbourhood of $E$, contained
in $O$.  On $Q$, $g=F$ is continuous and piecewise affine, hence
is a finite PL map.  The compact set $P\setminus\operatorname{Int}_P Q$
has closed image under $F$, disjoint from $y$.  A small target
neighbourhood $V_y\subset V$ avoids this image, so
$F^{-1}(V_y)\subset\operatorname{Int}_P Q$.  Thus
$F:F^{-1}(V_y)\to V_y$ is exactly the restriction of the finite
PL map $g|_Q$, with its entire fibre.  In the planar target chart,
take a small closed box $B$ about $y$ contained in that chart.
The union $S=g(Q)\cup B$ is a compact target polyhedron; shrink
$V_y$ further into the interior of $B$.  Triangulate the finite
PL map $g|_Q:Q\to S$ (which need not be surjective), and pass
to the canonical realizations of these finite complexes to obtain
(A5.13.1).  The same construction works in any finite-dimensional
PL chart.  All arguments in this step work in every elementary
extension.

\medskip\noindent\textbf{3. A punctured edge admits semilinear agreement.}
The input from Shiota is the unnumbered statement on printed p.33
of \cite{A5-Shiota}, proved on pp.33--34: if $K$ is a finite
simplicial complex and
$H=(H_1,H_2):|K|\to M^2$ is continuous definable with both
coordinates nonnegative, there are a definable isotopy $\zeta_t$
of $|K|$ preserving each simplex, a number $\epsilon>0$, and a
definable relative neighbourhood $N$ of
$(0,\epsilon]\times\{0\}$ in
$(0,\epsilon]\times[0,\epsilon]$, such that
$H\zeta_1$ on $(H\zeta_1)^{-1}(N)$ extends to a PL map
$L:|K|\to M^2$.  Shrinking $N$ makes it relatively open and
restricts to $0<s<\epsilon$, $0\leq t<\epsilon$.

Consider a continuous map from a compact definable set to a
two-dimensional polyhedral target, an interior point $a$ of a
target edge, and the case in which that edge has exactly two
incident triangles.  Affinely flatten the two triangles onto
opposite sides of an axis; in a neighbourhood of $a$ use
coordinates $(s,t)$ with the edge given by $t=0$.  For a sufficiently
small $b>0$, the closed right rectangle
$D=[a,a+b]\times[-b,b]$ lies in this chart.  Its full preimage
$X_D$ is compact.  Triangulate $X_D$ compatibly with the closed
subsets given by $t\circ f\geq0$ and $t\circ f\leq0$:
\[
 h_0:P=|K|\xrightarrow{\sim}X_D,
 \qquad P_+,P_-\subset P,
 \qquad P_0=P_+\cap P_- .
\]
These three subsets are subcomplexes.  Apply Shiota's statement to
\[
 H=(s\circ f\circ h_0-a,\ |t\circ f\circ h_0|):P\longrightarrow M^2.
\]
Because $\zeta_1$ preserves the subcomplexes, the homeomorphism
$h=h_0\zeta_1:P\to X_D$ preserves their interpretation as the two
sides.  Define a finite piecewise affine map $g:P\to M^2$ by
\[
 g(x)=
 \begin{cases}
 (a+L_1(x), L_2(x)),&x\in P_+\setminus P_0,\\
 (a+L_1(x),-L_2(x)),&x\in P_-\setminus P_0,\\
 (a+L_1(x),0),&x\in P_0.
 \end{cases}
 \tag{A5.13.4}
\]
It may fail to be continuous outside the neighbourhood in question;
Step~2 was formulated to allow exactly this possibility.
The signed inverse image of $N$ under $(s,t)\mapsto(s-a,|t|)$,
restricted to the interior of $D$, is an open set $V$ containing
$(a,a+\delta)\times\{0\}$ for some $\delta>0$.  On
$(f h)^{-1}(V)$, equation (A5.13.4) agrees with the signed map
$f h$.  Indeed, the unsigned coordinates agree with $L$ there,
the side is preserved, and $L_2=0$ on $P_0$ there.  Step~2 gives
a local finite PL model at every point of this deleted right
edge interval.  The left interval is obtained by using $a-s$.

Retain the stronger data, not merely the conclusion: a rectangle
$D$, a polyhedron $P$, a homeomorphism $h:P\to f^{-1}(D)$,
a finite piecewise affine map $g$, a target open set
$V\subset\operatorname{Int}D$, and a number $\delta>0$, such that
$V$ contains the specified deleted interval and $fh=g$ on
$(fh)^{-1}(V)$.  For each fixed choice of formulas and finite
piece data these are first-order requirements, of the kind
described in Step~1.

Fix one open target edge $e$.  At every $a\in e$, in every
elementary extension, the preceding construction supplies both
left and right certificates.  Applying Step~1 to pairs of such
formats yields finitely many formats covering all $a\in e$.
Let $D_e\subset e$ be the union of all deleted intervals certified
by all parameter instances of these finitely many formats.
This is a definable set: membership is an existential formula
with a finite disjunction.  For every $a\in e$ there is a
$\delta_a>0$ for which
\[
       (a-\delta_a,a+\delta_a)\cap e\setminus\{a\}\subset D_e.
 \tag{A5.13.5}
\]
The definable complement $e\setminus D_e$ has no interval, and
therefore is finite.  Every point of $D_e$ has a local finite PL
model by Step~2, in every elementary extension, because the
certificates asserting (A5.13.3) transfer.  Apply the local-model
consequence (A5.13.2) of Step~1 to the already definable set
$A=D_e$.  It proves definability of the maximal triviality locus
on $D_e$.  Its value on the remaining finite set $e\setminus D_e$
is also definable with parameters.  The desired locus on the
entire open edge is consequently definable.

\medskip\noindent\textbf{4. Reduction to the compact planar situation.}
Let $Y_{\rm man}$ denote the definable manifold-without-boundary
locus of $Y$.  Its definability was proved in A5.005.  The
components of local dimension zero or one are settled by A5.008.
Restrict to the union $T$ of its components of dimension two and
to $f^{-1}(T)\to T$.  This map remains definably proper.  The
set $T$ is locally compact and locally closed in its ambient
space.  One direct verification of the latter fact is to take
at each point a closed small ball inside a definable manifold
chart; its image is compact and is an ambient closed
neighbourhood in $T$.  Thus $T$ is open in its ambient closure.

Apply coordinatewise semialgebraic bounded homeomorphisms to
source and target, and let $\overline\Gamma$ be the closure of
the resulting graph inside a closed bounded box.  Put
$\overline T=\operatorname{cl}T$ in the bounded target coordinates.
Projection gives a continuous map of compact definable sets
\[
          \overline f:\overline\Gamma\longrightarrow\overline T.
\]
There are no added graph points over $T$.  In fact, at any
$y\in T$ take a compact target neighbourhood $C\subset T$ whose
interior contains $y$.  Properness makes its original preimage
compact; its transformed graph is compact and closed.  Any
point of $\overline\Gamma$ over $y$ is in its closure, since the
target coordinate of a convergent definable curve eventually
lies in the interior of $C$.  It therefore belongs to the
original graph.  This also proves
$\overline f^{-1}(T)\cong f^{-1}(T)$ over $T$.

Take a Hardt partition of $\overline f$, and a finite definable
triangulation of $\overline T$ compatible with that partition
and with $T$.  Its dimension is at most two.  Transfer the map
to this polyhedral target.  On every open two-simplex contained
in $T$, the full map is definably trivial by the chosen Hardt
partition, so all its points belong to the maximal open locus.
An open one-simplex contained in the two-manifold $T$ has exactly
two incident two-simplexes; its usual finite-polyhedral
neighbourhood is therefore the planar chart used in Step~3.
That step proves definability of the locus on every such edge.
There are finitely many remaining vertices.  Any subset of them
is definable with parameters.  The union over the finitely many
simplexes is the full maximal locus in $T$, and hence is
definable.  Joining the already settled components of dimensions
zero and one proves the theorem.
\end{proof}
